\documentclass[11pt]{article}
\usepackage[margin=1in]{geometry}
\usepackage{amsmath,amssymb}
\usepackage{enumitem}
\usepackage{graphicx}
\usepackage{booktabs}
\usepackage{xcolor}
\usepackage{titlesec}

\definecolor{garnet}{HTML}{73000A}
\definecolor{darkgray90}{HTML}{363636}

\titleformat{\section}{\Large\bfseries\color{garnet}}{}{0em}{}
\titleformat{\subsection}{\large\bfseries\color{darkgray90}}{}{0em}{}
\titleformat{\subsubsection}{\normalsize\bfseries\color{darkgray90}}{}{0em}{}

\setlength{\parindent}{0pt}
\setlength{\parskip}{6pt}

\title{\textbf{\color{garnet}CSCE 763 --- Digital Image Processing\\[4pt] Distance Measurement \& Mathematical Tools}}
\author{Study Notes}
\date{}

\begin{document}
\maketitle
\thispagestyle{empty}

\section{Distance Measurement}

A function $D(p, q)$ qualifies as a \textit{distance metric} if it satisfies three properties for pixels $p$, $q$, and $z$:

\begin{enumerate}[label=(\alph*)]
    \item \textbf{Non-negativity:} $D(p,q) \geq 0$, and $D(p,q) = 0$ if and only if $p = q$
    \item \textbf{Symmetry:} $D(p,q) = D(q,p)$
    \item \textbf{Triangle inequality:} $D(p,z) \leq D(p,q) + D(q,z)$
\end{enumerate}

Given pixel $p$ at coordinates $(x, y)$ and pixel $q$ at coordinates $(s, t)$, there are three standard distance measures.

\subsection{Euclidean Distance}

\begin{align}
    D_e(p,q) = \sqrt{(x-s)^2 + (y-t)^2}
\end{align}

This is the straight-line distance. Pixels at a given Euclidean distance from a center form a \textit{circle}.

\subsection{City-Block (D4) Distance}

\begin{align}
    D_4(p,q) = |x - s| + |y - t|
\end{align}

Movement is restricted to horizontal and vertical directions only, like navigating a Manhattan street grid. Pixels at a given D4 distance from a center form a \textit{diamond} shape. This distance corresponds to the $N_4$ neighborhood concept.

\subsection{Chessboard (D8 / Chebyshev) Distance}

\begin{align}
    D_8(p,q) = \max(|x - s|, |y - t|)
\end{align}

Movement is permitted in any of 8 directions, like a king in chess. Pixels at a given D8 distance from a center form a \textit{square}. This corresponds to the $N_8$ neighborhood.

\subsection{Sample Problem}

Consider pixels $p$ at the top-left corner and $q$ at the bottom-right corner of a 6-row, 3-column region. The row difference is 5 and the column difference is 1.

\begin{align}
    D_4 &= |5| + |1| = 6 \\
    D_8 &= \max(5, 1) = 5 \\
    D_e &= \sqrt{1 + 25} = \sqrt{26}
\end{align}

\textbf{Key distinction:} Distance is a direct measurement between two points. The \textit{length of a path} is the sum of distances between consecutive pixels along a connected path, which is always $\geq$ the direct distance.

\newpage
\section{Array vs.\ Matrix Operations}

This is a critical distinction in image processing.

\subsection{Array Operations (Element-wise)}

Each element in the result is computed from the corresponding elements in the same position. For multiplication:

\begin{align}
    \begin{bmatrix} a_{11} & a_{12} \\ a_{21} & a_{22} \end{bmatrix}
    \cdot
    \begin{bmatrix} b_{11} & b_{12} \\ b_{21} & b_{22} \end{bmatrix}
    =
    \begin{bmatrix} a_{11}b_{11} & a_{12}b_{12} \\ a_{21}b_{21} & a_{22}b_{22} \end{bmatrix}
\end{align}

Position $(i,j)$ in the output is simply the product of position $(i,j)$ from each input. This is what image arithmetic (averaging, subtraction, multiplication) uses --- pixel-by-pixel operations.

\subsection{Matrix Operations (Standard Linear Algebra)}

Uses dot products of rows and columns:

\begin{align}
    \begin{bmatrix} a_{11} & a_{12} \\ a_{21} & a_{22} \end{bmatrix}
    \times
    \begin{bmatrix} b_{11} & b_{12} \\ b_{21} & b_{22} \end{bmatrix}
    =
    \begin{bmatrix} a_{11}b_{11}+a_{12}b_{21} & a_{11}b_{12}+a_{12}b_{22} \\ a_{21}b_{11}+a_{22}b_{21} & a_{21}b_{12}+a_{22}b_{22} \end{bmatrix}
\end{align}

Matrix multiplication is used for \textit{geometric transformations} (affine transforms like rotation, scaling, translation, shear) and vector/matrix operations in color image processing.

\textbf{Bottom line:} Most pixel-level image processing uses \textit{array} (element-wise) operations. \textit{Matrix} multiplication is reserved for coordinate transformations and linear algebra tasks.

\section{Linear vs.\ Nonlinear Operations}

An operator $H$ is \textit{linear} if it satisfies the superposition property:

\begin{align}
    H\big[a_i f_i(x,y) + a_j f_j(x,y)\big] = a_i H\big[f_i(x,y)\big] + a_j H\big[f_j(x,y)\big]
\end{align}

This combines two properties:

\textbf{Additivity:} $H[f_1 + f_2] = H[f_1] + H[f_2]$. The response to a sum of inputs equals the sum of individual responses.

\textbf{Homogeneity:} $H[a \cdot f] = a \cdot H[f]$. Scaling the input scales the output by the same factor.

\textbf{Examples of linear operations:} image averaging, spatial convolution with a fixed kernel, Fourier transforms.

\textbf{Examples of nonlinear operations:} median filtering, thresholding, taking the max or min of pixel values. These break the superposition property because, for example, the median of a sum is not the sum of the medians.

Linearity matters because linear operations are mathematically tractable --- they can be analyzed in the frequency domain, their behavior can be predicted, and complex operations can be decomposed into simpler ones.

\newpage
\section{Applications of Image Arithmetic}

\subsection{Image Averaging --- Noise Reduction}

A noisy image is modeled as $g(x,y) = f(x,y) + \eta(x,y)$, where $f$ is the true image and $\eta$ is random noise. If $K$ noisy images of the same scene are captured and averaged:

\begin{align}
    \bar{g}(x,y) = \frac{1}{K}\sum_{i=1}^{K} g_i(x,y)
\end{align}

Then the expected value converges to the true image and the noise variance is reduced:

\begin{align}
    E\{\bar{g}(x,y)\} &= f(x,y) \\
    \sigma^2_{\bar{g}} &= \frac{1}{K}\sigma^2_{\eta}
\end{align}

Averaging $K$ images reduces noise by a factor of $K$. The assumption is that the noise is \textit{uncorrelated} between images and has \textit{zero mean}.

\subsection{Image Subtraction --- Enhancing Differences}

\begin{align}
    g(x,y) = f_1(x,y) - f_2(x,y)
\end{align}

There are two key applications. The first is \textbf{detecting changes}: subtracting two similar images highlights what changed between them. The second is \textbf{Digital Subtraction Angiography (DSA)}: a ``mask'' X-ray image is taken before injecting contrast dye, and a ``live'' image is taken after injection. Subtracting them cancels out all static anatomy (bones, tissue) and reveals only the blood vessels filled with dye.

\subsection{Image Multiplication --- ROI Masking}

\begin{align}
    g(x,y) = f(x,y) \cdot h(x,y)
\end{align}

A binary mask $h$ is created with 1s where the region of interest should be preserved and 0s elsewhere. Multiplying the mask with the image isolates the desired region. For example, a dental X-ray can be multiplied with a rectangular mask to isolate only the teeth with fillings.

\subsection{Image Division --- Shading Correction}

\begin{align}
    g(x,y) = f(x,y) \,/\, h(x,y)
\end{align}

If an image has non-uniform illumination (common in microscopy), the shading pattern $h(x,y)$ can be estimated and divided out. This corrects gradual brightness falloff across the image.

\subsection{Important Notes on Arithmetic Operations}

Images used in averaging and subtraction \textbf{must be registered} (spatially aligned). Misalignment will produce artifacts instead of meaningful results.

Output images should be \textbf{normalized to $[0, 255]$} using:
\begin{align}
    f_m &= f - \min(f) \\
    f_s &= K\left[\frac{f_m}{\max(f_m)}\right], \quad K = 255
\end{align}

\end{document}
